\clearpage
\VerdierPhysicalPageAnchor{0119}%
\begingroup
\small
\setlength{\abovedisplayskip}{4pt plus 2pt minus 2pt}%
\setlength{\belowdisplayskip}{4pt plus 2pt minus 2pt}%

\VerdierSemanticAnchor{ch2-1-2-8-splitting-complete}
$fgf=f$ かつ $gfg=g$。射は分裂をもつとき\emph{分裂可能}であるという。
分解可能な射は分裂可能である。逆は一般には成り立たない
（例：非単項 Dedekind 環上の自由加群の圏）。

\VerdierSemanticAnchor{ch2-1-2-8-decomposable-category}
加法圏は、そのすべての射が分解可能であるとき\emph{分解可能}であるという。
分解可能な圏はアーベル圏である。

\VerdierSemanticAnchor{ch2-1-2-8-triangle-decomposition}
$X\xrightarrow{f}Y$ を三角圏 $\mathcal D$ の射とし、

\VerdierSemanticAnchor{ch2-1-2-8-decomposition-diagram-repeat}
\begin{equation*}
  \begin{tikzcd}[ampersand replacement=\&,column sep=5.2em,row sep=3.4em]
    X \arrow[r,"f"] \arrow[d,"\wr"']
      \& Y \arrow[d,"\wr"] \\
    Z\oplus W
      \arrow[r,"{\left(\begin{smallmatrix}0&\operatorname{id}_W\\0&0\end{smallmatrix}\right)}"]
      \& W\oplus Z'
  \end{tikzcd}\,.
\end{equation*}
を $f$ の分解とする。$f$ を含む任意の完全三角は、三角

\VerdierSemanticAnchor{ch2-1-2-8-decomposed-triangle}
\begin{equation*}
  \begin{tikzcd}[ampersand replacement=\&,column sep=2.0em]
    Z\oplus W
      \arrow[r,"{\left(\begin{smallmatrix}0&\operatorname{id}_W\\0&0\end{smallmatrix}\right)}"]
      \& W\oplus Z'
      \arrow[r,"{\left(\begin{smallmatrix}0&\operatorname{id}_{Z'}\\0&0\end{smallmatrix}\right)}"]
      \& Z'\oplus Z[1]
      \arrow[r,"{\left(\begin{smallmatrix}0&\operatorname{id}_{Z[1]}\\0&0\end{smallmatrix}\right)}"]
      \& Z[1]\oplus W[1]
  \end{tikzcd}\,.
\end{equation*}
と同型である。

\VerdierSemanticAnchor{ch2-1-2-8-decomposed-triangle-derivation}
これは (1.2.4)、(1.2.5)、および公理 TRI、TRII から直ちに従う。

\VerdierSemanticAnchor{ch2-1-2-9}
\medskip
\noindent\textbf{命題 1.2.9.}\quad
\emph{$\mathcal D$ を三角圏とし、$f\colon X\longrightarrow Y$ を
$\mathcal D$ の射とする。射 $f$ に関する次の性質を考える。}

\VerdierSemanticAnchor{ch2-1-2-9-properties}
\begin{enumerate}
\renewcommand{\labelenumi}{\roman{enumi})}
\item \emph{射 $f$ は核をもつ。}
\item \emph{射 $f$ は余核をもつ。}
\item \emph{射 $f$ は分解可能である。}
\item \emph{射 $f$ は分裂可能である。}
\end{enumerate}

\noindent\emph{このとき、次の含意が成り立つ。}

\VerdierSemanticAnchor{ch2-1-2-9-implications}
\[
  \mathrm{i)}\Longleftrightarrow\mathrm{ii)}
  \Longleftrightarrow\mathrm{iii)}\Longrightarrow\mathrm{iv)}\,.
\]

\VerdierSemanticAnchor{ch2-1-2-9-countable-equivalence}
\emph{$\mathcal D$ において可算積または可算直和が表現可能ならば、
性質 \textnormal{i)}、\textnormal{ii)}、\textnormal{iii)}、
\textnormal{iv)} は同値である。}

\VerdierSemanticAnchor{ch2-1-2-9-proof-opening}
まず、$\mathrm{iii)}\Longrightarrow\mathrm{i)}$、
$\mathrm{iii)}\Longrightarrow\mathrm{ii)}$、
$\mathrm{iii)}\Longrightarrow\mathrm{iv)}$ は明らかである。含意
$\mathrm{i)}\Longrightarrow\mathrm{iii)}$ は、$\mathcal D$ の反対圏に
おける含意
\endgroup
